\paragrafos{Ο γενικός όρος α·ν: προβλέπω οποιονδήποτε όρο}{Α.4.2, «Κανονικότητες της μορφής α·ν με α ρητό»}

\begin{anagke}
Η Πολυξένη αγοράζει από το περίπτερο πακέτα με κάρτες για το άλμπουμ της. Κάθε πακέτο έχει $5$ κάρτες. Με $1$ πακέτο έχει $5$ κάρτες, με $2$ πακέτα $10$, με $3$ πακέτα $15$. «Πόσες κάρτες θα έχω αν αγοράσω $100$ πακέτα;» ρωτάει η Αγγελική. Ο Στέλιος αρχίζει να γράφει: $5,\ 10,\ 15,\ 20,\ 25,\ \ldots$ και σταματάει στο $60$ να παραπονιέται: «Δεν θα γράφω εκατό αριθμούς!».

Ο Ηλίας φτιάχνει πίνακα με τη σειρά και τον αριθμό των καρτών και βλέπει κάτι: $5=5\cdot1$, $10=5\cdot2$, $15=5\cdot3$, $20=5\cdot4$. «Κάθε αριθμός είναι το $5$ επί τη θέση του! Ο εκατοστός είναι ο $5\cdot100$». Η Ρήνη ρωτάει το αντίστροφο: «Αν έχω $225$ κάρτες, πόσα πακέτα αγόρασα;». Η Αγγελική προσθέτει: «Και αν ένα ρομποτάκι προχωρά $\frac{3}{4}$ του μέτρου κάθε δευτερόλεπτο; Ο κανόνας λειτουργεί και τότε;».

Η κυρία Αλεξάνδρα χαμογελά: «Αν βρείτε έναν τύπο που δίνει τον όρο για οποιαδήποτε θέση, μπορείτε να τον υπολογίσετε χωρίς να γράψετε τους προηγούμενους». Σε αυτή την παράγραφο μαθαίνουμε τι είναι ο γενικός όρος μιας κανονικότητας, τον βρίσκουμε όταν είναι της μορφής $\alpha\cdot\nu$ (και όταν ο $\alpha$ είναι ρητός, όχι μόνο φυσικός) και τον χρησιμοποιούμε για να βρίσκουμε όρους και θέσεις.
\end{anagke}

\begin{lexi}{ακολουθία, ν-οστός όρος, γενικός όρος}{ακολουθία, ν-οστός όρος, γενικός όρος}
Η \lexiword{ακολουθία} είναι αυτό που «ακολουθεί»: ένα στοιχείο μετά το άλλο, με τη σειρά τους. Ο \lexiword{ν-οστός όρος} είναι ο όρος που βρίσκεται στη θέση $\nu$, όπως λέμε «πέμπτος», «εκατοστός». Ο \lexiword{γενικός όρος} είναι «γενικός» επειδή ισχύει για κάθε θέση: ένας τύπος που περιγράφει όλους τους όρους μαζί.
\end{lexi}

\begin{metafrasi}
Ο όρος που βρίσκεται στη θέση $\nu$. & $\nu$-οστός όρος\\
Πέντε επί τη θέση του όρου. & $5\cdot\nu$ ή $5\nu$\\
Ο $10$ος όρος της κανονικότητας $5\nu$. & $5\cdot10=50$\\
Ποιος όρος ισούται με $225$; & $5\nu=225$, άρα $\nu=225:5=45$\\
Κάθε όρος είναι $\alpha$ φορές η θέση του. & $\alpha\cdot\nu$\\
\end{metafrasi}

\tmima{Η θεωρία}

\tmimasmall{Ακολουθίες και κανονικότητες}

Οι αριθμοί $5,\ 10,\ 15,\ 20,\ 25,\dots$ αποτελούν μια \textbf{ακολουθία}. Ακολουθία είναι μια διαδοχή στοιχείων, που ονομάζονται \textbf{όροι}, και έχουν μια ορισμένη σειρά: $1$ος όρος, $2$ος όρος, $3$ος όρος κ.ο.κ.

\begin{center}
\renewcommand{\arraystretch}{1.3}
\begin{tabular}{rcccccc}
$5,$ & $10,$ & $15,$ & $20,$ & $25,$ & $\dots$\\
$\uparrow$ & $\uparrow$ & $\uparrow$ & $\uparrow$ & $\uparrow$ & \\
{\small $1$ος όρος} & {\small $2$ος όρος} & {\small $3$ος όρος} & {\small $4$ος όρος} & {\small $5$ος όρος} &
\end{tabular}
\end{center}

Αν κάθε όρος μιας ακολουθίας προκύπτει με έναν συγκεκριμένο κανόνα, τότε οι όροι της αποτελούν μια \textbf{κανονικότητα}.

\begin{orismos}{ν-οστός όρος ή γενικός όρος}
Ο όρος που βρίσκεται στη θέση $\nu$ ($\nu=1,\,2,\,3,\dots$) λέγεται \textbf{ν-οστός όρος} ή \textbf{γενικός όρος} της ακολουθίας.
\end{orismos}

\tmimasmall{Η κανονικότητα $\alpha\cdot\nu$}

Στην κανονικότητα $5,\ 10,\ 15,\ 20,\ 25,\dots$ κάθε όρος είναι κατά $5$ μεγαλύτερος από τον προηγούμενο (κανόνας: $+5$). Παρατηρούμε ότι:
\begin{itemize}[leftmargin=1.2em,itemsep=0.3ex,topsep=2pt]
\item $1$ος όρος ($\nu=1$): $5=5\cdot1$
\item $2$ος όρος ($\nu=2$): $10=5\cdot2$
\item $3$ος όρος ($\nu=3$): $15=5\cdot3$
\item $4$ος όρος ($\nu=4$): $20=5\cdot4$
\item $\nu$-οστός όρος: $5\cdot\nu$
\end{itemize}

Ο γενικός όρος είναι λοιπόν $5\cdot\nu$ ή, πιο σύντομα, $5\nu$. Αν γνωρίζουμε τον γενικό όρο, μπορούμε να βρούμε όποιον όρο θέλουμε:
\begin{itemize}[leftmargin=1.2em,itemsep=0.3ex,topsep=2pt]
\item Ο $10$ος όρος είναι $5\cdot10=50$ και ο $23$ος είναι $5\cdot23=115$.
\item Για να βρούμε ποιος όρος ισούται με $225$, γράφουμε $5\nu=225$, άρα $\nu=225:5=45$. Ο $45$ος όρος ισούται με $225$.
\end{itemize}

\begin{center}
\begin{tikzpicture}[x=1.1cm,y=0.13cm]
  \draw[-{Stealth},thick] (0,0)--(5.6,0) node[right,font=\small]{$\nu$};
  \draw[-{Stealth},thick] (0,0)--(0,28) node[above,font=\small]{όρος};
  \foreach \x in {1,...,5}{\draw (\x,0.6)--(\x,-0.6) node[below,font=\scriptsize]{$\x$};}
  \foreach \y in {5,10,15,20,25}{\draw (0.07,\y)--(-0.07,\y) node[left,font=\scriptsize]{$\y$};}
  \draw[dashed,gray] (0,0)--(5.4,27);
  \foreach \x in {1,...,5}{\pgfmathtruncatemacro{\yy}{5*\x}\fill[gblue] (\x,\yy) circle (2.4pt);}
\end{tikzpicture}
\end{center}

Τα σημεία $(\nu,\ \text{όρος})$ της κανονικότητας $5\nu$ βρίσκονται σε μία ευθεία γραμμή, που περνά από την αρχή των αξόνων.

\begin{kanonas}{κανονικότητες της μορφής $\alpha\cdot\nu$}
Αν ο $1$ος όρος μιας κανονικότητας είναι ο $\alpha$ και κάθε όρος είναι κατά $\alpha$ μεγαλύτερος από τον προηγούμενο, τότε ο γενικός όρος είναι $\alpha\cdot\nu$. Ο $\alpha$ είναι ταυτόχρονα ο πρώτος όρος και το «βήμα». Αν γνωρίζω έναν όρο και τη θέση του, βρίσκω τον $\alpha$ διαιρώντας: $\alpha=\text{όρος}:\nu$.
\end{kanonas}

Ο $\alpha$ μπορεί να είναι οποιοσδήποτε ρητός: φυσικός ($7\nu$), αρνητικός ($-3\nu$), δεκαδικός ($0{,}5\nu$) ή κλάσμα ($\frac{3}{4}\nu$). Η σειρά $\nu$ είναι πάντα φυσικός αριθμός ($1,\,2,\,3,\dots$), γι' αυτό ένας αριθμός είναι όρος της κανονικότητας $\alpha\nu$ μόνο όταν το πηλίκο «αριθμός $:\ \alpha$» είναι φυσικός.

\tmimasmall{Γεωμετρικά μοτίβα}

Ένα ισόπλευρο τρίγωνο με πλευρά από $\nu$ οδοντογλυφίδες έχει συνολικά $3\nu$ οδοντογλυφίδες. Μια πυραμίδα με βάση πολύγωνο $\nu$ πλευρών έχει $\nu$ ακμές στη βάση και άλλες $\nu$ που ενώνουν την κορυφή με τη βάση: $2\nu$ ακμές. Σε κάθε περίπτωση ο γενικός όρος μας λέει πόσα στοιχεία χρειάζομαι για οποιοδήποτε σχήμα της σειράς.

\begin{prosoxi}
\begin{itemize}[leftmargin=1.2em,itemsep=0.3ex,topsep=2pt]
\item Το $7\nu$ δεν είναι το $7+\nu$: το πρώτο δίνει $7,\,14,\,21,\dots$ και το δεύτερο $8,\,9,\,10,\dots$
\item Το $\nu$ είναι η \emph{θέση} και όχι ο όρος. Ο $225$ ως όρος της $5\nu$ είναι ο $45$ος όρος, όχι ο $225$ος.
\item Για να βρω τη θέση διαιρώ: $\nu=\text{όρος}:\alpha$. Δεν πολλαπλασιάζω.
\item Αν το πηλίκο «όρος $:\ \alpha$» δεν βγαίνει φυσικός, ο αριθμός δεν είναι όρος της κανονικότητας: ο $100$ δεν είναι όρος της $7\nu$, γιατί $100:7$ δεν είναι φυσικός.
\item Δεν είναι κάθε προσθετική κανονικότητα της μορφής $\alpha\nu$: η $4,\,7,\,10,\,13,\dots$ έχει βήμα $3$, αλλά ο πρώτος όρος είναι $4$, όχι $3$. Ο γενικός της όρος είναι άλλος ($3\nu+1$).
\end{itemize}
\end{prosoxi}

\tmima{Λυμένα παραδείγματα}

\paradeigma{Βρίσκω τον γενικό όρο}
\tixero Η κανονικότητα $6,\ 12,\ 18,\ 24,\dots$\par
\tithelo Ο γενικός όρος, ο $7$ος και ο $20$ός όρος.\par
\possindeo Κάθε όρος είναι κατά $6$ μεγαλύτερος από τον προηγούμενο και ο πρώτος είναι $6$: $6=6\cdot1$, $12=6\cdot2$, $18=6\cdot3$.\par
\apantisi Γενικός όρος $6\nu$. Ο $7$ος είναι $6\cdot7=42$ και ο $20$ός $6\cdot20=120$.\par
\elegxos Ο $4$ος όρος από τον τύπο είναι $6\cdot4=24$, όπως στη σειρά.

\paradeigma{Από τη θέση στον όρο και αντίστροφα}
\tixero Ο γενικός όρος είναι $7\nu$.\par
\tithelo Ο $20$ός όρος. Ποιος όρος ισούται με $511$; Υπάρχει όρος ίσος με $100$;\par
\possindeo Ο $20$ός όρος: $\nu=20$. Για το $511$ λύνω $7\nu=511$, άρα $\nu=511:7=73$. Για το $100$: $100:7=14{,}28\ldots$ που δεν είναι φυσικός.\par
\apantisi $7\cdot20=140$. Ο $73$ος όρος ισούται με $511$. Όχι, δεν υπάρχει όρος ίσος με $100$.\par
\elegxos $7\cdot73=511$. Και τα πολλαπλάσια του $7$ δίπλα στο $100$ είναι τα $98$ και $105$.

\paradeigma{Αρνητικός συντελεστής}
\tixero Η κανονικότητα $-4,\ -8,\ -12,\ -16,\dots$\par
\tithelo Ο γενικός όρος, ο $25$ος όρος και η θέση του όρου $-96$.\par
\possindeo Κάθε όρος είναι κατά $4$ μικρότερος (βήμα $-4$) και ο πρώτος είναι $-4$: $-4=(-4)\cdot1$, $-8=(-4)\cdot2$. Για τη θέση: $(-4)\nu=-96$, άρα $\nu=(-96):(-4)=24$.\par
\apantisi Γενικός όρος $-4\nu$. Ο $25$ος είναι $-100$. Ο $-96$ είναι ο $24$ος όρος.\par
\elegxos $(-4)\cdot24=-96$.

\paradeigma{Δεκαδικός συντελεστής}
\tixero Η κανονικότητα $0{,}5,\ 1,\ 1{,}5,\ 2,\ \ldots$\par
\tithelo Ο γενικός όρος, ο $30$ός όρος και η θέση του όρου $12{,}5$.\par
\possindeo Ο πρώτος όρος είναι $0{,}5$ και κάθε όρος αυξάνεται κατά $0{,}5$: $0{,}5\nu$. Ο $30$ός όρος: $0{,}5\cdot30=15$. Για το $12{,}5$: $0{,}5\nu=12{,}5$, άρα $\nu=12{,}5:0{,}5=25$.\par
\apantisi Γενικός όρος $0{,}5\nu$, ο $30$ός είναι $15$ και ο $12{,}5$ είναι ο $25$ος όρος.\par
\elegxos Το $0{,}5\nu$ σημαίνει «το μισό του $\nu$»: το μισό του $30$ είναι $15$ και το μισό του $25$ είναι $12{,}5$.

\paradeigma{Κλασματικός συντελεστής}
\tixero Η κανονικότητα $\frac{3}{4},\ \frac{3}{2},\ \frac{9}{4},\ 3,\ \ldots$\par
\tithelo Ο γενικός όρος, ο $8$ος όρος και η θέση του όρου $15$.\par
\possindeo Ο πρώτος όρος είναι $\frac{3}{4}$ και κάθε όρος αυξάνεται κατά $\frac{3}{4}$: $\frac{3}{4}\nu$. Ο $8$ος: $\frac{3}{4}\cdot8=6$. Για το $15$: $\frac{3}{4}\nu=15$, άρα $\nu=15:\frac{3}{4}=15\cdot\frac{4}{3}=20$.\par
\apantisi Γενικός όρος $\frac{3}{4}\nu$, ο $8$ος είναι $6$ και ο $15$ είναι ο $20$ός όρος.\par
\elegxos $\frac{3}{4}\cdot20=15$.

\paradeigma{Πυραμίδες}
\tixero Μια πυραμίδα με βάση τρίγωνο έχει $6$ ακμές, με τετράπλευρη βάση $8$, με πενταγωνική $10$.\par
\tithelo Ο γενικός όρος και η πυραμίδα που έχει $50$ ακμές.\par
\possindeo Οι ακμές είναι $6,\,8,\,10,\dots$ για βάσεις με $3,\,4,\,5,\dots$ πλευρές, δηλαδή $6=2\cdot3$, $8=2\cdot4$, $10=2\cdot5$. Άρα $2\nu$ ακμές για βάση με $\nu$ πλευρές. Για τις $50$ ακμές: $2\nu=50$, $\nu=25$.\par
\apantisi Γενικός όρος $2\nu$ και η βάση έχει $25$ πλευρές.\par
\elegxos Η βάση με $25$ πλευρές χρειάζεται $25$ ακμές στη βάση και $25$ ακμές προς την κορυφή: $50$.

\paradeigma{Πίνακας με άγνωστους όρους}
\tixero Οι όροι της κανονικότητας με γενικό όρο $\alpha\nu$ είναι: $1{,}2,\ 2{,}4,\ x,\ 4{,}8,\ y,\ 7{,}2$.\par
\tithelo Οι $x$ και $y$ και ο $100$ός όρος.\par
\possindeo Ο πρώτος όρος είναι $1{,}2$, άρα $\alpha=1{,}2$ και ο γενικός όρος είναι $1{,}2\nu$. Ο $x$ είναι ο $3$ος και ο $y$ ο $5$ος όρος. Ο $6$ος είναι $1{,}2\cdot6=7{,}2$, όπως δίνεται.\par
\apantisi $x=1{,}2\cdot3=3{,}6$, $y=1{,}2\cdot5=6$ και ο $100$ός όρος $1{,}2\cdot100=120$.\par
\elegxos Η διαφορά διαδοχικών όρων είναι $1{,}2$: $2{,}4+1{,}2=3{,}6$ και $3{,}6+1{,}2=4{,}8$.

\tmima{Ασκήσεις}

\zoni{Το κατάλαβα}{γενικός όρος, ν-οστός όρος, θέση}

\ex Συμπλήρωσε τον πίνακα με τη βοήθεια του κανόνα. Σειρά $\nu$: $1,\ 2,\ 3,\ 4,\ 5,\ 6$ και όροι: $9,\ 18,\ 27,\ 36,\ \_\_,\ \_\_$. Ποιος από τους τύπους α) $9+\nu$ \ β) $9\nu$ \ γ) $9+9\nu$ είναι ο γενικός όρος; Ποιος από τους α) $4+\nu$ \ β) $4\nu$ \ γ) $4+4\nu$ \ δ) $8\nu-4$ είναι ο γενικός όρος της σειράς $4,\ 8,\ 12,\ 16,\dots$;

\ex Βρες τους $5$ πρώτους όρους της κανονικότητας με γενικό όρο: α) $8\nu$ \quad β) $11\nu$ \quad γ) $-3\nu$ \quad δ) $0{,}5\nu$.

\ex Συμπλήρωσε τον πίνακα και βρες τον γενικό όρο: α) $6,\ 12,\ 18,\ 24,\ \_\_,\ \_\_$ \quad β) $-2,\ -4,\ -6,\ -8,\ \_\_,\ \_\_$ \quad γ) $0{,}25,\ 0{,}5,\ 0{,}75,\ 1,\ \_\_,\ \_\_$.

\ex Η κανονικότητα έχει γενικό όρο $9\nu$. Βρες τον $10$ο, τον $25$ο και τον $100$ό όρο.

\ex Η κανονικότητα έχει γενικό όρο $6\nu$. Ποιος όρος ισούται με $72$; Με $240$; Με $1\,500$;

\ex Είναι όρος της κανονικότητας $7\nu$ ο αριθμός: α) $91$ \quad β) $150$ \quad γ) $700$ \quad δ) $5$; Αν ναι, ποια είναι η θέση του;

\ex Σωστό ή λάθος; α) Ο γενικός όρος της σειράς $5,\,10,\,15,\dots$ είναι $5+\nu$. \quad β) Ο $4$ος όρος της $3\nu$ είναι $12$. \quad γ) Στην κανονικότητα $2\nu$ ο $1$ος όρος είναι $2$. \quad δ) Ο γενικός όρος της $-3,\,-6,\,-9,\dots$ είναι $-3\nu$. \quad ε) Στον τύπο $4\nu$ το $\nu$ μπορεί να είναι $0{,}5$. \quad στ) Αν ο γενικός όρος είναι $\alpha\nu$, ο $1$ος όρος είναι ο $\alpha$. \quad ζ) Ο $9$ος όρος της $2\nu$ είναι $18$.

\zoni{Το χρησιμοποιώ}{οδοντογλυφίδες, πυραμίδες, παγωτά, ρομπότ}

\ex \textbf{Οι οδοντογλυφίδες.} Ο Βασίλης φτιάχνει ισόπλευρα τρίγωνα με πλευρά $\nu$ οδοντογλυφίδες. α) Πόσες οδοντογλυφίδες χρειάζεται για τρίγωνο με πλευρά $7$; β) Πόση πλευρά έχει το τρίγωνο που χρειάζεται $120$ οδοντογλυφίδες; γ) Υπάρχει τρίγωνο της σειράς που χρειάζεται ακριβώς $100$ οδοντογλυφίδες;

\ex \textbf{Οι πυραμίδες.} Μια πυραμίδα με βάση πολύγωνο $\nu$ πλευρών έχει $\nu$ ακμές στη βάση και $\nu$ ακμές προς την κορυφή. Συμπλήρωσε τον πίνακα για βάση με $3,\,4,\,5,\,6$ πλευρές. Βρες τις ακμές πυραμίδας με βάση $8$-γωνο, $9$-γωνο και $20$-γωνο και τον γενικό όρο.

\ex \textbf{Τα παγωτά.} Ο Νίκος αγοράζει παγωτά των $1{,}5$ € το ένα. Πόσο κοστίζουν $\nu$ παγωτά; Πόσο κοστίζουν $12$; Πόσα παγωτά αγοράζει με $27$ €;

\ex \textbf{Το ρομποτάκι.} Το ρομποτάκι της Ζωής προχωρά $\frac{3}{4}$ του μέτρου κάθε δευτερόλεπτο. Πόσα μέτρα έχει διανύσει μετά από $12$ δευτερόλεπτα; Σε πόσα δευτερόλεπτα θα διανύσει $15$ μέτρα;

\ex \textbf{Το θερμόμετρο.} Σε μια παγωμένη νύχτα η θερμοκρασία ξεκινά από $0$° και πέφτει $1{,}5$° κάθε ώρα. Γράψε τον γενικό όρο της θερμοκρασίας μετά από $\nu$ ώρες. Πόση θα είναι μετά από $8$ ώρες; Πότε θα φτάσει τους $-27$°;

\zoni{Το σκέφτομαι}{ποιος όρος; ποιο βήμα;}

\ex Ο Στέλιος γράφει ότι ο γενικός όρος της $6,\ 12,\ 18,\ 24,\dots$ είναι $6+6\nu$. Πώς μπορείς να ελέγξεις με τον $1$ο όρο ότι κάνει λάθος; Ποιος είναι ο σωστός τύπος;

\ex Μια κανονικότητα έχει γενικό όρο $\alpha\nu$ και ο $4$ος όρος της είναι $10$. Βρες τον $\alpha$, τους δύο πρώτους όρους και τον $50$ό όρο.

\ex Βρες όλους τους φυσικούς $\alpha$ για τους οποίους ο αριθμός $30$ είναι όρος της κανονικότητας $\alpha\nu$.

\ex Η Ζωή γράφει τη σειρά $3,\ 6,\ 9,\ 12,\dots$ και ο Ηλίας τη σειρά $4,\ 7,\ 10,\ 13,\dots$ Και οι δύο έχουν βήμα $3$. Ποια από τις δύο έχει γενικό όρο $3\nu$; Γιατί η άλλη όχι;

\zoni{Για το σπίτι}{πίνακες και καταστάσεις}

\ex Βρες δύο καθημερινές καταστάσεις που περιγράφονται από τον γενικό όρο $12\nu$ και από τον $0{,}5\nu$ (π.χ. τα αυγά σε $\nu$ δωδεκάδες).

\ex Για την κανονικότητα με γενικό όρο $2{,}5\nu$, φτιάξε πίνακα με τους $8$ πρώτους όρους και σημείωσε τα σημεία σε σύστημα αξόνων. Τι παρατηρείς για τη θέση των σημείων;
